\documentclass[aspectratio=169,10pt,english]{beamer}

\input{../../../_preambulo_beamer.tex}
\input{../../../_comandos_beamer.tex}

\selectlanguage{english}

\title{MA1001B - Statistical Modeling for Decision Making}
\subtitle{Lesson 06: Counting Techniques I}
\author{}
\institute{Tecnol\'ogico de Monterrey}
\date{}

\begin{document}

\begin{frame}[plain]
  \vspace{1.2cm}
  {\Large\bfseries MA1001B - Statistical Modeling for Decision Making\par}
  \vspace{0.7cm}
  {\large Lesson 06: Counting Techniques I\par}
  \vspace{0.7cm}
  {\color{TecRojo}\rule{\textwidth}{1.2pt}}
  \vspace{0.8cm}

  MA1001B Faculty\\[0.3cm]
  Term\\[0.3cm]
  Tecnol\'ogico de Monterrey
\end{frame}

\begin{frame}{The Counting Question}
  \begin{alertblock}{Guiding question}
    How can we determine the size of a finite sample space without listing every
    outcome, and how can we detect overcounting?
  \end{alertblock}

  \vspace{0.5em}
  By the end of this lesson, you can:
  \begin{itemize}
    \item apply the product rule to successive choice stages;
    \item distinguish ordered assignments from unordered subsets;
    \item verify formulas through exhaustive enumeration;
    \item diagnose repeated counting when restrictions interact.
  \end{itemize}

  \vfill
  \scriptsize All facilities, roles, people, and restrictions are synthetic.
\end{frame}

\begin{frame}{Step 1: Product Rule for Complete Configurations}
  \begin{columns}[T]
    \begin{column}{0.40\textwidth}
      \small
      Choose one element from each stage:
      \[
        N=\prod_{i=1}^{m}n_i.
      \]

      \begin{block}{Synthetic review configuration}
        3 fulfillment centers\\
        4 workflow stages\\
        2 evidence types
      \end{block}

      \[
        3\times4\times2=24
      \]

      Exhaustive Python enumeration also returns 24 tuples. Each center
      contributes 8 configurations.
    \end{column}
    \begin{column}{0.60\textwidth}
      \centering
      \includegraphics[width=\textwidth]{../figures/counting_01_product_rule.png}
    \end{column}
  \end{columns}
\end{frame}

\begin{frame}{Step 2: Permutations for Distinct Roles}
  \begin{columns}[T]
    \begin{column}{0.40\textwidth}
      \small
      Assign Lead, Validation, and Reporting roles from six analysts.

      \[
        P(n,k)=\frac{n!}{(n-k)!}
      \]
      \[
        P(6,3)=6\times5\times4=120
      \]

      \begin{exampleblock}{Why order matters}
        The same three analysts in different roles create a different
        assignment. Python generates exactly 120 ordered tuples.
      \end{exampleblock}
    \end{column}
    \begin{column}{0.60\textwidth}
      \centering
      \includegraphics[width=\textwidth]{../figures/counting_02_permutations.png}
    \end{column}
  \end{columns}
\end{frame}

\begin{frame}{Step 3: Combinations for an Unranked Panel}
  \begin{columns}[T]
    \begin{column}{0.40\textwidth}
      \small
      Select three analysts for a peer panel with no internal roles.

      \[
        \binom{n}{k}=\frac{n!}{k!(n-k)!}
      \]
      \[
        \binom{6}{3}=20
      \]

      Each three-person subset has $3!=6$ possible orderings:
      \[
        \frac{P(6,3)}{3!}=\frac{120}{6}=20.
      \]
    \end{column}
    \begin{column}{0.60\textwidth}
      \centering
      \includegraphics[width=\textwidth]{../figures/counting_03_combinations.png}
    \end{column}
  \end{columns}
\end{frame}

\begin{frame}{The Outcome Determines the Counting Method}
  \small
  \begin{center}
    \begin{tabular}{p{0.27\textwidth}p{0.26\textwidth}p{0.34\textwidth}}
      \toprule
      \textbf{Question} & \textbf{Method} & \textbf{Lesson 06 result} \\
      \midrule
      One choice at every stage & Product rule & $3\times4\times2=24$ configurations \\
      Distinct positions assigned & Permutation & $P(6,3)=120$ role assignments \\
      Unranked subset selected & Combination & $\binom{6}{3}=20$ panels \\
      \bottomrule
    \end{tabular}
  \end{center}

  \vspace{0.8em}
  \begin{alertblock}{Decision test}
    Describe one final outcome before choosing a formula. If exchanging two
    selected elements changes that outcome, order matters.
  \end{alertblock}

  A formula can return an integer and still answer the wrong counting question.
\end{frame}

\begin{frame}{Step 4: The Overcounting Limit}
  \begin{columns}[T]
    \begin{column}{0.43\textwidth}
      \small
      Choose four people from five operations analysts and four data specialists.
      Require both disciplines and exclude the pair O1 with D1.

      \vspace{0.4em}
      \begin{alertblock}{Naive expression}
        $5\times4\times\binom{7}{2}=420$\\
        The entire universe has only $\binom{9}{4}=126$ groups.
      \end{alertblock}

      \[
        126-5-1-21=99
      \]
      Exhaustive enumeration confirms 99 valid groups.
    \end{column}
    \begin{column}{0.57\textwidth}
      \centering
      \includegraphics[width=\textwidth]{../figures/counting_04_restrictions.png}
    \end{column}
  \end{columns}
\end{frame}

\begin{frame}{Concept Check}
  \begin{enumerate}
    \item A review code uses 4 centers, 3 stages, and 2 evidence types. How many
      complete configurations exist?
    \item How many ordered assignments fill two distinct roles from five people?
    \item How many unranked two-person panels can be selected from five people?
    \item Why does a result of 420 immediately reveal an error when the full
      universe contains only 126 unique groups?
  \end{enumerate}

  \vspace{0.5em}
  \begin{block}{Connection to Lesson 07}
    The next lesson extends counting to repeated categories and multinomial
    allocations, where several elements share the same type.
  \end{block}

  \vfill
  \scriptsize Source: OpenStax, \textit{Introductory Business Statistics 2e},
  Chapter 3, Section 3.1. CC BY-NC-SA 4.0.
\end{frame}

\appendix



\end{document}
